% 恒星光谱（综述）
% license CCBYSA3
% type Wiki

本文根据 CC-BY-SA 协议转载翻译自维基百科\href{https://en.wikipedia.org/wiki/Stellar_classification}{相关文章}。

在天文学中，恒星分类是根据恒星的光谱特性对恒星进行的分类。来自恒星的电磁辐射通过棱镜或衍射光栅被分解成光谱，呈现出类似彩虹的连续颜色，其间夹杂着光谱线。每一条光谱线都对应一种特定的化学元素或分子，而谱线的强度则反映该元素的丰度。不同光谱线强度的变化主要由恒星光球层的温度决定，尽管在某些情况下也确实存在真实的元素丰度差异。恒星的光谱型是一种简短的代码，主要概括了其电离状态，从而为光球层温度提供一个客观的度量。

目前，大多数恒星采用摩根–基南（Morgan–Keenan，MK）分类系统进行分类，该系统使用字母 O、B、A、F、G、K、M，按从最热（O 型）到最冷（M 型）的顺序排列。每一个字母类别又进一步用一个数字细分，其中 0 表示最热，9 表示最冷（例如 A8、A9、F0、F1 就构成了一个由热到冷的序列）。此外，该序列还扩展出 W、S、C 三个类别，用于那些不符合经典体系的恒星。一些恒星遗迹或质量明显偏离恒星范围的天体也被赋予了字母标记：D 表示白矮星，L、T、Y 表示棕矮星（以及部分系外行星）。

在 MK 系统中，还会在光谱型后加上一个用罗马数字表示的光度级。光度级依据恒星光谱中某些吸收线的宽度来判定，这些线宽随恒星大气密度变化，因此可用于区分巨星与矮星。其中，光度级 0 或 Ia+ 表示特超巨星，I 表示超巨星，II 表示亮巨星，III 表示普通巨星，IV 表示亚巨星，V 表示主序星，sd（或 VI）表示亚矮星，而 D（或 VII）表示白矮星。太阳的完整光谱分类为 G2V，这表明太阳是一颗主序星，其表面温度约为 5800 K。
\subsection{传统颜色描述}
传统的恒星颜色描述只考虑恒星光谱的峰值位置。然而实际上，恒星会在整个电磁波谱范围内辐射。由于所有光谱颜色叠加后呈现为白色，人眼实际看到的恒星颜色要比传统颜色描述所暗示的浅得多。这种“明亮度”的特性说明，将恒星简单地赋予某种光谱颜色可能具有误导性。在排除暗光条件下的颜色对比效应后，在一般观测条件中，并不存在真正的绿色、青色、靛色或紫色恒星。像太阳这样的“黄色”矮星实际上呈现为白色；“红矮星”呈现为较深的黄色或橙色；而“棕矮星”并不会真的呈现棕色，而是假想中在近距离观察时可能显得暗红色，或呈灰色/黑色。
\subsection{现代分类}
现代恒星分类体系被称为摩根–基南（Morgan–Keenan，MK）分类系统。在该系统中，每一颗恒星都会被赋予一个光谱型（源自较早的哈佛光谱分类体系，该体系本身并不包含光度级$^{[1]}$），并再配以一个用罗马数字表示的光度级（如下所述），二者共同构成恒星的完整光谱类型。

其他现代恒星分类体系（例如 UBV 系统）则基于颜色指数，即通过测量两个或多个波段星等之间的差值来进行分类$^{[2]}$。这些数值通常以诸如 U−V、B−V这样的形式给出，表示通过两种标准滤光片（例如紫外、蓝光和可见光）所测得的颜色差异。
\subsubsection{哈佛光谱分类}
哈佛体系是一种一维的恒星分类方案，由天文学家安妮·江普·坎农（Annie Jump Cannon）提出，她对德雷珀（Draper）早期按字母顺序排列的分类体系进行了重新排序与简化（见“历史”部分）。在该体系中，恒星根据其光谱特征被划分为若干组，每一组用一个字母表示，并可辅以数字细分。

主序星的表面温度大致分布在 2000–50 000 K 之间，而演化程度更高的恒星——尤其是新近形成的白矮星——其表面温度则可超过 100 000 K$^{[3]}$。从物理意义上看，这些光谱类别反映的是恒星大气的温度，并通常按由热到冷的顺序排列。
\begin{figure}[ht]
\centering
\includegraphics[width=14.25cm]{./figures/1615d7b930f17d79.png}
\caption{} \label{fig_Stcl_1}
\end{figure}
用于记忆光谱型字母由热到冷排列顺序的传统助记口诀是“Oh, Be A Fine Guy/Girl: Kiss Me!”（“哦，做个好小伙/好姑娘：吻我吧！”）$^{[12]}$。
尽管在天文学课程和相关组织举办的活动中，人们提出过许多替代性的助记口诀，但这一传统口诀至今仍然是最为流行的$^{[13][14]}$。

光谱类型 O 到 M（以及后文讨论的其他更为专门的类型）都会再用阿拉伯数字 0–9进行细分，其中 0 表示该类型中最热的恒星。例如，A0 表示 A 型中最热的恒星，而 A9 则表示最冷的 A 型恒星。允许使用小数；例如，恒星 Mu Normae 的光谱型被定为 O9.7 $^{[15]}$。太阳的光谱型为 G2$^{[16]}$。

尽管哈佛光谱分类实际上反映的是恒星的表面温度或光球温度（更准确地说是其有效温度），但这一点在该体系建立之初并未被完全理解。不过，到第一张赫罗图（Hertzsprung–Russell diagram）于 1914 年前后提出时，人们已经普遍怀疑这一对应关系是成立的$^{[17]}$。在 1920 年代，印度物理学家梅格纳德·萨哈（Meghnad Saha）在物理化学中分子解离理论的基础上，发展出了电离理论并将其推广到原子的电离问题。他首先将该理论应用于太阳色球层，随后又应用于恒星光谱$^{[18]}$。

随后，哈佛天文学家塞西莉亚·佩恩（Cecilia Payne）证明O–B–A–F–G–K–M这一光谱序列本质上就是一个温度序列$^{[19]}$。由于该分类序列的建立早于人们对其“温度本质”的认识，将一条光谱归入某一具体亚型（如 B3 或 A7）在很大程度上依赖于对恒星光谱中吸收特征强度的（带有一定主观性的）估计。因此，这些亚型并不是在数学意义上等间隔划分的。
\subsubsection{摩根–基南（Morgan–Keenan）分类}
耶基斯光谱分类，也称为 MK 分类，或 Morgan–Keenan（亦称 MKK，Morgan–Keenan–Kellman）系统$^{[20][21]}$，是一种由威廉·威尔逊·摩根（William Wilson Morgan）、菲利普·C·基南（Philip C. Keenan）和伊迪丝·凯尔曼（Edith Kellman）于 1943 年在耶基斯天文台提出的恒星光谱分类体系$^{[22]}$。这是一个二维分类体系（温度与光度），其依据是对恒星温度和表面重力敏感的光谱线；而表面重力又与恒星的光度密切相关（相比之下，哈佛分类仅基于表面温度）。在 1953 年对标准星表和分类标准进行若干修订之后，该体系被正式命名为摩根–基南分类（MK）$^{[23]}$，并一直沿用至今。

表面重力更高、密度更大的恒星，其光谱线会表现出更显著的压力展宽。由于巨星的半径远大于质量相近的矮星，其表面重力（以及由此产生的压力）要低得多。因此，光谱中的差异可以被解释为光度效应，从而仅通过分析光谱本身就能够确定恒星的光度级。

目前区分出了若干不同的光度级，具体列于下表中$^{[24]}$。
\begin{figure}[ht]
\centering
\includegraphics[width=14.25cm]{./figures/73151fcf04a39c9e.png}
\caption{} \label{fig_Stcl_2}
\end{figure}
允许存在边缘情况；例如，一颗恒星可能既可被归为超巨星也可被归为亮巨星，或者介于亚巨星与主序星之间。在这些情况下，会在两个光度级之间使用两种特殊符号：
\begin{itemize}
\item 斜杠（/）：表示恒星可能属于其中任意一个光度级。
\item 连字符（-）：表示恒星处在两个光度级之间的过渡状态。
\end{itemize}
例如，被标记为 A3-4III/IV 的恒星，表示其光谱型介于 A3 与 A4 之间，同时其光度级介于巨星（III）与亚巨星（IV）之间，或可能属于其中之一。

此外，还曾使用亚矮星（sub-dwarf）光度级：

VI 表示亚矮星（其光度略低于主序星）。

名义上的光度级 VII（有时甚至更高的罗马数字）如今已很少用于白矮星或“热亚矮星”的分类，因为主序星和巨星所采用的温度字母体系已不再适用于白矮星。

在某些情况下，字母 a 和 b 也会用于非超巨星的光度级中；例如，一颗光度略低于典型巨星的恒星可被标为 IIIb，而 IIIa 则表示其光度略高于典型巨星$^{[34]}$。

一小部分 V 型极端恒星在 He II λ4686 吸收谱线处表现出很强的吸收特征，因此被赋予 Vz 的分类标记。一个典型例子是恒星 HD 93129 B$^{[35]}$。
\subsubsection{光谱异常}
还可以在光谱型之后附加小写字母形式的补充标记，用以指出光谱中的某些特殊或异常特征$^{[36]}$。
\begin{table}[ht]
\centering
\caption{请输入表格标题}\label{tab_Stcl1}
\begin{tabular}{|c|c|}
\hline
\textbf{代码} & \textbf{恒星的光谱异常} \\
\hline
: & 不确定的光谱值$^{[24]}$\\
\hline
... & 存在未描述的光谱异常\\
\hline
! & 特殊异常\\
\hline
comp & 复合光谱[37]\\
\hline
e & 存在发射线[37]\\
\hline
[e] & 存在“禁戒”发射线\\
\hline
er & “反向”发射线中心比边缘弱\\
\hline
eq & 具有P Cygni轮廓的发射线\\
\hline
f & N III 和 He II 发射$^{[24]}$\\
\hline
f & NIV 4058Å线比NIII 4634Å、4640Å和4642Å线更强$^{[38]}$\\
\hline
f+ & 除了N III线外，还发射Si IV 4089Å和4116Å线$^{[38]}$\\
\hline
f? & C III 4647–4650–4652Å发射线的强度与N III线相当$^{[39]}$\\
\hline
(f) & N III发射，He II缺失或吸收较弱\\
\hline
(f+) & $^{[40]}$\\
\hline
((f)) & 显示强烈的He II吸收，伴随微弱的N III发射$^{[41]}$\\
\hline
((f*)) & $^{[40]}$\\
\hline
h & WR恒星以氢发射线为特征。$^{[42]}$\\
\hline
ha & WR恒星的氢线在吸收和发射中均可见。$^{[42]}$\\
\hline
He wk & 弱氦线\\
\hline
k & 带有星际吸收特征的光谱\\
\hline
m & 增强的金属特征$^{[37]}$\\
\hline
n & 因自旋引起的宽泛（“模糊”）吸收$^{[37]}$\\
\hline
nn & 非常宽泛的吸收特征$^{[24]}$\\
\hline
neb & 星云的光谱混入其中$^{[37]}$\\
\hline
p & 未指定的特殊性，奇特恒星.\\
\hline
pq & 奇特光谱，类似于新星的光谱\\
\hline
q & P型仙后座轮廓\\
\hline
s & 窄（“锐利”）吸收线$^{[37]}$\\
\hline
ss & 极窄线\\
\hline
sh & 壳层星特征$^{[37]}$\\
\hline
变量 & 可变光谱特征$^{[37]}$（有时简称为“v”）\\
\hline
波长 & 弱线$^{[37]}$（也称“w”和“wk”）\\
\hline
元素符号 & 指定元素的异常强烈的光谱线$^{[37]}$\\
\hline
z & 表明在468.6 纳米处存在异常强烈的电离氦线$^{[35]}$\\
\hline
\end{tabular}
\end{table}
例如，59 Cygni 被列为光谱型 B1.5Vnne$^{[43]}$，这表示其光谱的总体分类为 B1.5V，同时还具有非常宽的吸收线以及某些发射线。
\subsection{历史}
哈佛光谱分类中字母顺序之所以显得有些“反常”，其原因在于历史演变：该体系起源于更早的塞基（Secchi）分类，并随着人们对恒星物理认识的不断加深而逐步修订和重组。
\subsubsection{塞基分类}
在 19 世纪 60—70 年代，恒星光谱学的先驱安杰洛·塞基（Angelo Secchi） 创建了塞基光谱分类体系，用于对观测到的恒星光谱进行分类。到 1866 年，他已经提出了三类恒星光谱，如下表所示$^{[44][45][46]}$。

在 19 世纪末（1890 年代后期），这一分类体系逐渐被哈佛光谱分类所取代；后者正是本文其余部分所讨论的现代恒星光谱分类基础$^{[47][48][49]}$。
\begin{table}[ht]
\centering
\caption{请输入表格标题}\label{tab_Stcl2}
\begin{tabular}{|c|c|}
\hline
\textbf{等级编号} & \textbf{塞奇等级描述}\\
\hline
塞奇等级I & 白蓝相间的恒星，具有宽阔而强烈的氢线，例如织女星和牛郎星。这包括现代A型和早期F型。\\
\hline
塞奇等级I（猎户座亚型）& $<b$一种属于Secchi I类的亚型，其光谱中以窄线取代了宽带，例如参宿四和大角星。按现代术语来说，这对应于早期B型恒星\\
\hline
塞奇等级II & 黄色恒星——氢线较弱，但金属谱线明显，例如太阳, 大角星和五车二。这包括现代的G型和K型恒星，以及晚期F型恒星。\\
\hline
塞奇等级III & Orange to red stars with complex band spectra, such as Betelgeuse and Antares.
这对应于现代的M类。\\
\hline
塞奇四等 & In 1868, he discovered carbon stars, which he put into a distinct group:$^{[50]}$带有显著碳谱带和谱线的红色恒星，对应于现代的C类和S类。\\
\hline
塞奇五等 & 1877年，他增设了第五等：$^{[51]}$发射线星，例如仙后座γ星和谢利阿克星, 属于现代的Be型。1891年，爱德华·查尔斯·皮克林提出，V类应对应于现代的O型（当时还包括沃尔夫-拉叶星）。以及行星状星云中的恒星。\\
\hline
\end{tabular}
\end{table}
用于塞基（Secchi）光谱分类的罗马数字，不应与耶基斯（Yerkes）光度分类以及拟议的中子星分类中所使用的、彼此毫不相关的罗马数字混淆。
\subsubsection{德雷珀体系}
\begin{figure}[ht]
\centering
\includegraphics[width=14.25cm]{./figures/a92bc94a88248741.png}
\caption{} \label{fig_Stcl_3}
\end{figure}
在丈夫去世后，玛丽·安娜·德雷珀（Mary Anna Draper）开始资助哈佛底片库（Harvard Plate Stacks）的建立，以及在哈佛学院天文台对这些天文底片的研究。天文台台长爱德华·C·皮克林（Edward C. Pickering）随后雇用了多位开创性的女性天文学家，她们后来被统称为“哈佛计算员”（Harvard Computers）。尽管最初设想她们会研究多个不同的天文学课题，但这项工作的一个早期重要成果，是《亨利·德雷珀恒星光谱纪念目录》（The Henry Draper Memorial Catalogue of Stellar Spectra）的第一版，该目录于 1890 年首次出版。

威廉敏娜·弗莱明（Williamina Fleming）完成了该目录第一版中大部分光谱的分类工作，她被认为共分类了 一万余颗恒星，并发现了 10 颗新星以及 200 多颗变星。在哈佛计算员，尤其是威廉敏娜·弗莱明的帮助下，亨利·德雷珀目录的最初版本被设计出来，用以取代由安杰洛·塞基（Angelo Secchi）建立的罗马数字光谱分类体系。

该目录采用了一种新的分类方案：将此前使用的塞基分类（I 至 V 类）进一步细分为更具体的类别，并用 A 到 P 的字母表示；此外，还使用字母 Q 来标记那些无法归入其他任何类别的恒星。$^{[53][54]}$弗莱明与皮克林合作，根据氢谱线强度的不同区分出了 17 个光谱类别。氢谱线强度的变化会导致恒星辐射波长的差异，从而造成颜色外观的变化。在这一体系中，A 类光谱表现出最强的氢吸收线，而 O 类光谱几乎不显示可见谱线。随着字母顺序向后推进，氢吸收线的强度逐渐减弱。后来，安妮·跳·坎农（Annie Jump Cannon）和安东尼娅·莫里（Antonia Maury）对该体系进行了修订，最终形成了哈佛光谱分类体系。$^{[56][58]}$
\subsubsection{旧哈佛体系（1897）}
1897 年，哈佛的另一位天文学家安东尼娅·莫里将塞基 I 类中的“猎户型”亚类置于塞基 I 类其余部分之前，从而在分类顺序上使现代的 B 型位于现代的 A 型之前。她是第一位这样做的人，尽管她并未使用字母标记的光谱类型，而是采用了一个由 I–XXII 共 22 个编号组成的分类序列。$^{[59][60]}$
\begin{table}[ht]
\centering
\caption{1897年哈佛系统的概要$^{[61]}$}\label{tab_Stcl3}
\begin{tabular}{|c|c|}
\hline
\textbf{分组} & \textbf{摘要}\\
\hline
I−V & 包括了从I组到V组氢吸收线强度逐渐增强的‘猎户座型’恒星\\
\hline
VI & 充当了‘猎户座型’与Secchi I型之间的中间类型\\
\hline
VII−XI & 是塞基1型恒星，其氢吸收线强度从第VII组至第XI组逐渐减弱\\
\hline
XIII−XVI & 包括塞基2型恒星，其氢吸收线强度逐渐减弱，而太阳型金属线则逐渐增强\\
\hline
XVII−XX & $<b$包括光谱线逐渐增多的塞奇3型恒星\\
\hline
XXI & 包括第4型天狼星\\
\hline
二十二 & 包括沃尔夫-拉叶星\\
\hline
\end{tabular}
\end{table}
由于最初的 22 个罗马数字分组仍不足以涵盖恒星光谱中出现的全部变化，人们又引入了三种附加划分，以进一步细化差异：在原有编号后添加小写字母，用来区分谱线外观的相对特征；这些谱线被定义为：
\begin{itemize}
\item $a$：平均宽度
\item $b$：模糊
\item $c$：锐利
\end{itemize}
安东尼娅·莫里**于 **1897 年**发表了她自己的恒星分类目录，题为
《**用 11 英寸德雷珀望远镜拍摄的亮星光谱：亨利·德雷珀纪念项目的一部分**》（*Spectra of Bright Stars Photographed with the 11 inch Draper Telescope as Part of the Henry Draper Memorial*）。
该目录收录了 **4,800 张照片**以及莫里本人对 **681 颗北天亮星**的分析。这是**首次由女性署名的天文台正式出版物**。
\subsubsection{现行哈佛体系（1912）}
**1901 年**，**安妮·跳·坎农**重新采用了字母分类法，但她舍弃了除 **O、B、A、F、G、K、M、N** 之外的所有字母，并按这一顺序使用；此外，她还保留了 **P**（用于行星状星云）和 **Q**（用于某些异常光谱）。她还引入了诸如 **B5A**（介于 B 型与 A 型之间的恒星）、**F2G**（从 F 型向 G 型过渡五分之一）的混合标记方式，依此类推。

最终，到 **1912 年**，坎农将 **B、A、B5A、F2G** 等类型统一改写为 **B0、A0、B5、F2** 等形式。这一体系基本上就是**现代哈佛光谱分类系统**。该系统建立在对**摄影底片上恒星光谱的系统分析**之上，这种方法能够将恒星发出的光转换为可读的光谱。
\subsubsection{威尔逊山分类}
一种称为**威尔逊山体系**的**光度分类法**曾被用于区分不同光度的恒星。即便在今天，这种标记方式仍偶尔可在现代光谱中见到。其主要符号包括：

\begin{itemize}
\item sd：次矮星（subdwarf）
\item d：矮星（dwarf）
\item sg：次巨星（subgiant）
\item g**：巨星（giant）
\item c**：超巨星（supergiant）
\end{itemize}

\subsection{光谱型}

“**光谱型**”（Spectral type）一词在此指恒星光谱分类；若指小行星的光谱分类，则另见“小行星光谱类型”。

恒星分类体系本质上是一种**分类学（taxonomy）体系**，以**模式恒星**为基础，类似于生物学中对物种的分类：每一个类别及其子类别都由一颗或多颗**标准恒星**来定义，并配有相应的**判别特征描述**。
